%!TEX program = uplatex
\documentclass{ujarticle}

\setlength{\topmargin}{-1.4cm}
%\setlength{\paperwidth}{14cm}
\setlength{\oddsidemargin}{-30Q}
\setlength{\evensidemargin}{0cm}

\setlength{\columnsep}{2.5zw}

\usepackage{amsmath}
\usepackage{amssymb}
\usepackage{amsthm}
\usepackage{eucal}


%%%%%%%%%%%%%%%%%%%%%%%%%%%%%%%%
%%%%%% 以下の項目を適宜，修正して
%%%%%% 利用してください．
%%%%%%
%%%%%%%%%%%%%%%%%%%%%%%%%%%%%%%%
\newcommand{\TITLE}%
{擬測地的な距離空間への群作用に関する Švarc-Milnor Lemma について}%卒論タイトルを書く
\newcommand{\eTITLE}%
{A Study on the Švarc-Milnor Lemma for Group Actions on Quasi-geodesic Metric Spaces}%%英語タイトル
\newcommand{\STNO}%
{2264138}%学籍番号
\newcommand{\NAME}%
{杉山 和輝}%氏名
\newcommand{\ADVR}%
{野崎 雄太 准教授}%指導教員氏名

%%%% 以下は変更しないでください
\newcommand{\DATE}%
{(2026年2月6日)}%〆切
\newcommand{\NENDO}%
{2025年度}%
%%%%%%%%%%%%%%%%%%%%%%%%%%%%%%%%%%

\setlength{\textwidth}{52zw}
\setlength{\textheight}{25cm}


\makeatletter

\def\ps@lrheadf{\ps@empty
  \def\@evenhead{\normalfont%\footnotesize
      \leftmark{}{}\hfil \rightmark{}{}}%
  \let\@oddhead\@evenhead
  \def\@oddfoot{}
  \let\@evenfoot\@oddfoot
  \let\@mkboth\markboth}
\makeatother


\pagestyle{lrheadf}

\markboth{\NENDO \hskip1zw 
数物・電子情報系学科 数理科学EP 
\hskip1zw
卒業研究 \hskip1zw 概要}{}



\begin{document}
%\fontsize{9pt}{14pt}\selectfont
\fontsize{10pt}{14pt}\selectfont


\hfil
\textbf{\Large \TITLE }

\vskip5Q
\hfil
{\large \eTITLE }


\vskip5Q
\hfil 
{\large \STNO \hskip1zw \NAME
\hskip 2zw 指導教員： \ADVR}

\vskip20Q


\noindent
\textbf{１．研究背景・動機}

群を理解することを目的とした分野の一つに，幾何学的群論が挙げられる．幾何学的群論とは，群から幾何学的対象を構成する，あるいは他の幾何学的対象への群作用をみることで，群の構造や代数的・幾何学的性質を考察する分野である．群から幾何学的対象を構成する代表的な例としては，群の元を頂点とし，生成系の元を演算することで移る頂点間を辺で結ぶことで得られるCayleyグラフが挙げられる．生成系によってグラフの形は変わるものの，生成系込みでの群の構造を可視化するこができる．またCayleyグラフは群作用する空間の代表例といえる．一般に群作用とは，ある集合の自己同型写像として群が振る舞うことであり，群の代数的構造を空間の幾何学的な構造として表現するものである．群を直接調べるのではなく，どのような空間にどのように作用しているかを観察することで，空間の幾何学的な特徴や性質を通して群を調べることを可能としている。グラフの自己同型はサイクルなどのグラフの構造を保つため，Cayleyグラフの構造を群作用を通して理解することもできる．幾何学的群論では群作用に着目することが多い．群の二点間の距離を，Cayleyグラフの一辺を長さ1として見ることで得られる距離を群に導入することで，群を距離空間として捉えることができる．距離を入れた群の，他の距離空間への群作用を考えることで，群と空間の類似性を議論することができるのだが，Cayleyグラフの構造は選んだ生成系に依存してしまうため，群に導入した距離関数も生成系に依存してしまう．そこで，一定誤差のズレを許容する``大雑把な視点''を導入し局所的な違いを無視することで，大局的な部分だけを見比べ，群の本質的な構造を抽出することができ，Cayleyグラフの生成系に依存してしまう問題を解決できる．

Švarc-Milnor Lemma とは，「群がよい距離空間によい群作用をしているとき，その群の有限生成系が分かり，群の構造と距離空間の構造が``大雑把に''等しい」という主張であり，その内容から「幾何学的群論の基本補題」とも言われている．例えば，群 $\mathbb{Z}^2$ が距離空間 $\mathbb{R}^2$ に加法で群作用をしているとき，この補題から群である $\mathbb{Z}^2$ と距離空間 $\mathbb{R}^2$ が大雑把に等しい構造を持っていることが分かる．よく知られている距離空間へのうまい作用さえ見つけることができれば，群の構造を理解する足掛かりとなる．この補題は1955年に Švarc の論文~[1] にて用いられ，1968年に Milnor~[4] が Švarc とは独立して本補題を見つけ応用した．本論文では Löh~[3] の議論をもとに，Švarc-Milnor Lemma の証明と，その補題をいいかえた主張の証明を与える．

\vskip9Q
\noindent
\textbf{２．主結果 (Švarc-Milnor Lemma)}

群 $G$ が，よい性質を持った距離空間 $X$ に群作用をしており，この群作用が次の3条件を持っているとき，群 $G$ は有限生成であり，群 $G$ と距離空間 $X$ は大雑把に等しい構造を持っている：(1) 距離を保つ群作用である．(2) 直径が有限のある部分集合 $B\subset X$ がとれ， $G$ の作用で $X$ 全体を覆える．(3) $B$ を少し大きくした集合 $B^\prime$ と，$B^\prime$ を群作用で移した集合 $g\cdot B^\prime$が互いに素でないような $g\in G$ が有限である．

またこの補題を用いり，コンパクトなリーマン多様体の普遍被覆が，多様体自身の基本群と大雑把に等しい構造を持っていることを，先行研究にならい詳細な部分まで証明した．

\vskip9Q
\noindent
\textbf{３．意義・証明のアイデアや方法}

群 $G$ の元 $g$ に対し，距離空間 $X$ 上の点 $x$ と群作用した点 $g\cdot x$ の二点を結ぶ測地線の像から，$g$ を生成する元を選び取ることができる．また，群作用の条件から，その生成する元が有限個であることが示せる．

群 $G$ から 距離空間 $X$ への写像が特定の条件を満たせばよいのだが，それは $G$ 上に導入した距離空間の定義や群作用の特徴から導くことができる．

\vskip9Q
\noindent
\textbf{４．今後の課題}

本研究では補題の証明と，コンパクトなリーマン多様体という具体的な幾何構造を持つ空間に適用した例をまとめた．今後は他の具体的な幾何構造をもつ空間 (負曲率をもつ空間など) に適用した例を考え，研究・理解を深めたい．

\vskip9Q
\noindent
\textbf{参考文献}

{\small

\noindent
[1] A.S. Švarc, \textit{A volume invariant of coverings}, Doklady Akademii Nauk SSSR, vol.105, 1955, pp.32--34.

\noindent
[2] Bridson, Haefliger, \textit{Metric Spaces of N-Positive Curvature}, Spiiinger, 1999.

\noindent
[3] Clara Löh, \textit{Geometric Group Theory}, Springer, 2018.

\noindent
[4] J. Milnor, \textit{A note on curvature and fundamental group}, Journal of Differential Geometry, vol.2, No.1, 1968, pp.1--7.

\noindent
[5] V.A. Efremovich, \textit{The proximity geometry of Riemannian manifolds}, Uspekhi Matematicheskikh Nauk 8, No.5, 1953, pp.189--191.

\noindent
[6] 正井 秀俊，群と幾何をみる，日本評論社，2023．

}

\end{document}